\chapter{Command Cheat-Sheet}
\label{ch:cheatsheet}

Every command this manual actually uses, grouped by tool.

\section{Node / Next.js}
\begin{lstlisting}[style=olkcodebash, language=bash]
npm install          # install dependencies
npm run dev          # start the dev server (localhost:3000)
npm run build        # production build
npm run start        # run a production build locally
npm run lint         # ESLint
\end{lstlisting}

\section{Supabase CLI (via npx — nothing installed globally)}
\begin{lstlisting}[style=olkcodebash, language=bash]
npx supabase init                       # one-time: create supabase/config.toml
npx supabase start                      # start the local Docker stack
npx supabase stop                       # stop it (data persists in a volume)
npx supabase status                     # reprint local credentials
npx supabase db reset                   # wipe + replay every migration, local only
npx supabase migration new <name>       # scaffold a correctly-timestamped migration
npx supabase db dump --local --schema public > supabase/schema.sql
                                         # regenerate the schema.sql snapshot

# Hosted project (production) -- one-time link per machine (Chapter 4, S:link):
npx supabase login
npx supabase link --project-ref <ref>

# Hosted project (production) -- every push, once linked:
npx supabase migration list             # compare local vs. hosted migration history
npx supabase db push --dry-run          # preview which migrations would be applied
npx supabase db push                    # apply pending migrations to the hosted project
npx supabase db query "<sql>" --linked  # read-only query against the hosted project
\end{lstlisting}

\section{npm script wrappers (Chapter~\ref{ch:local-env} -- prefer these day-to-day)}
\begin{lstlisting}[style=olkcodebash, language=bash]
npm run env:local        # point .env.local at the local Docker stack
npm run env:remote       # point .env.local at the hosted project
                          # (restart `npm run dev` after either)

npm run db:reset         # = npx supabase db reset (schema only, wipes local data)
npm run db:reset-full    # db reset, then db:pull-data -- schema + real data, one command
npm run db:pull-data     # copy public/auth data from the hosted project into local
npm run db:push          # = npx supabase db push
\end{lstlisting}

\section{Docker (rarely needed directly — the Supabase CLI drives this for you)}
\begin{lstlisting}[style=olkcodebash, language=bash]
docker ps --filter "name=_olk"          # see which OLK containers are running
docker exec supabase_db_olk psql -U postgres -d postgres -c "\dt public.*"
                                         # list tables directly against the local DB
docker info                             # confirm the daemon is reachable at all
\end{lstlisting}

\section{Verifying the environment (Chapter~\ref{ch:local-env})}
\begin{lstlisting}[style=olkcodebash, language=bash]
docker --version
docker info           # confirms the daemon is reachable, not just the CLI
groups                # confirm your user is in the `docker` group on Linux
\end{lstlisting}

\section{Compiling this manual}
\begin{lstlisting}[style=olkcodebash, language=bash]
cd Guide
pdflatex main.tex && pdflatex main.tex   # run twice: once for content, once for
                                          # the table of contents and cross-references
# or, if latexmk is available:
latexmk -pdf main.tex
\end{lstlisting}

\begin{olktip}
Every \code{Chapter~??} you might see in a single-pass compile is not an error — it means the cross-reference hasn't been resolved yet because \LaTeX{} needs a second pass to know where every labelled chapter landed. Always compile twice (or use \code{latexmk}, which does this automatically) before judging the output.
\end{olktip}
